% ED26-SCI-JCGM2008-001; source-specific reconstruction and mathematical clarification.
\chapter{Measurement Uncertainty}
\label{sec:p15-measurement-uncertainty}

A measurement supports an inference about a specified quantity through a
measurement procedure, a model and the information used to evaluate its
inputs. An uncertainty calculation describes that information under the
model. It does not certify that the model includes every relevant effect.

This treatment follows Sections 2--8, Annex D and Annex E.3 of
JCGM 100:2008(E), with the end-gauge and electrical examples of H.1--H.2
\cite{JCGM1002008EMeasurement}. The 2026 amendment's qualification concerning
nonlinear output estimates is incorporated where it applies
\cite{JCGM100Amendment2026Nonlinearity}. Mathematical identities below have
explicit probability assumptions. Recommendations about specifying,
evaluating and reporting measurements are attributed metrological practice.
The treatment is not a claim of conformity of an actual measurement to a
current standard.

\section{Specify what is being measured}
\label{sec:p15-measurement-target}

\begin{hbdefinition}{def:p15-measurement-target}{Measurand, realization and result}
A \emph{measurand} is the particular quantity intended to be measured,
specified by its relevant object, conditions and interpretation. A
\emph{realization} is the quantity actually subjected to the procedure.
A \emph{measurement result} supplies an estimated value, its units and the
associated uncertainty information.

A specification may require a location, temperature, pressure, support
condition or time. Its necessary detail depends on the accuracy sought.
Uncertainty due to an incomplete specification is distinct from variation in
repeated observations of one fixed realization.
\end{hbdefinition}

For a sheet thickness specified at a reference temperature, a micrometer
observes a particular location under an applied pressure and an actual
temperature. Calibration, temperature and compression corrections connect
that observation to the target. Variation across locations does not disappear
when the same point is measured repeatedly. Either the location belongs in the
specification or the interpretation of the spatial variation belongs in the
uncertainty assessment \cite[Annex D.1--D.5]{JCGM1002008EMeasurement}.
Likewise, the guide's steel-bar example requires temperature and pressure
information for micrometre accuracy that may be immaterial for a millimetre
comparison \cite[3.1.3]{JCGM1002008EMeasurement}.

\begin{hbdefinition}{def:p15-measurement-error}{Error, correction and uncertainty}
Relative to a specified numerical reference value \(v\), the error of a
reported value \(y\) is \(y-v\). If \(v\) is unknown, the actual error is
unknown. A correction compensates for an estimated influence; for an
estimated additive error \(b\), the corresponding correction is \(-b\).

\emph{Standard uncertainty} expresses an assigned or estimated dispersion as
a standard deviation. It is not the realized error, the value of a correction
or a guaranteed error bound. A correction can be nonzero with small
uncertainty, or zero with substantial uncertainty.
\end{hbdefinition}

The guide assumes that recognized significant systematic effects have been
corrected and that their uncertain corrections remain in the model.
Calibration against a reference does not make the reference exact.
A small uncertainty also cannot protect a conclusion against an overlooked
effect or a blunder in transcription \cite[3.2--3.4]{JCGM1002008EMeasurement}.
A hypothetical error equal to zero and a large assigned uncertainty are
compatible: the actual error and the information available about it are
different objects.

The distinction in the guide's Figures D.1--D.2 can be expressed without
identifying an unobservable error with an uncertainty bar. The observed mean
is shifted by the applied correction; uncertainty in both contributes to
uncertainty of the corrected result. An unknown residual error locates that
result relative to the target value, but is not an additional observed
correction. Uncertainty from an incompletely specified target is an additional
component only when it has not already entered the model
\cite[D.6, Figs.~D.1--D.2]{JCGM1002008EMeasurement}.

\section{The measurement model and its probability interpretation}
\label{sec:p15-measurement-model}

\begin{hbdefinition}{def:p15-measurement-model}{Input and output uncertainty model}
Choose units and numerical coordinates for \(d\geq1\) input quantities.
Let \(X=(X_1,\ldots,X_d)^{\mathsf T}\) be a random vector representing
uncertainty about those inputs, with finite second moments, mean
\(m=EX\) and covariance matrix
\[
 \Sigma_{ij}=E[(X_i-m_i)(X_j-m_j)].
\]
A measurable function \(f\) specifies the numerical output \(Y=f(X)\).
The input estimates are \(x=(x_1,\ldots,x_d)\); their assessed covariance
matrix is \(U=(u_{ij})\), with \(u_{ii}=u_i^2\) and \(u_i\geq0\).
The guide writes \(u_i=u(x_i)\) and \(u_{ij}=u(x_i,x_j)\).

The matrix \(U\) must be symmetric and positive semidefinite. When
\(u_i u_j>0\), write \(r_{ij}=u_{ij}/(u_i u_j)\). Correlation is not defined
by this quotient if a denominator is zero; in a valid covariance matrix
\(u_i=0\) implies \(u_{ij}=0\) for every \(j\).
\end{hbdefinition}

A mean and variance do not determine a probability distribution. Even all
moments need additional determinacy assumptions; distinct lognormal-family
laws with the same moments are exhibited in
\cref{thm:p15-criteria-lognormal-family}. Moment calculations below assert
only the moments whose hypotheses establish them.

The physical quantity need not fluctuate because a probability distribution
represents uncertainty about it. Conversely, a distribution of repeated
readings describes sampling variability. The intended interpretation must be
stated before these objects are combined. The guide adopts a common
probabilistic treatment of both information sources
\cite[4.1.1, 4.1.6, E.3]{JCGM1002008EMeasurement}.
This is an interpretive choice about a model, not a proof that a particular
assigned distribution is correct.

The model includes uncertain calibration values, corrections, influence
quantities and reference data. It may be an algorithm rather than a closed
formula. A discrepancy against check standards or independent determinations
can require revising the model rather than enlarging a numerical uncertainty
without explanation \cite[3.4.1--3.4.2, 4.1.2]{JCGM1002008EMeasurement}.

\begin{hbexample}{def:p15-measurement-power}{A model with an uncertain correction}
For voltage \(V\), reference resistance \(R_0>0\), temperature coefficient
\(\alpha\) and temperatures \(t,t_0\), the guide uses
\[
 P=\frac{V^2}{R_0[1+\alpha(t-t_0)]},
 \qquad 1+\alpha(t-t_0)>0.
\]
Its sensitivities are
\[
 P_V=\frac{2P}{V},\quad P_{R_0}=-\frac P{R_0},\quad
 P_\alpha=-\frac{P(t-t_0)}{1+\alpha(t-t_0)},\quad
 P_t=-\frac{P\alpha}{1+\alpha(t-t_0)}.
\]
The first quotient is used only for \(V\ne0\); direct differentiation gives
\(P_V=0\) at zero. Nonuniform temperature or nonlinear resistance may require
a different model. These formulas specify the conditional calculation, not
the adequacy of the circuit assumptions
\cite[4.1.1--4.1.2, 5.1.3]{JCGM1002008EMeasurement}.
\end{hbexample}

\section{Evaluating input uncertainty}
\label{sec:p15-measurement-input-evaluation}

\begin{hbdefinition}{def:p15-measurement-type-ab}{Type A and Type B evaluation}
A \emph{Type A evaluation} uses statistical analysis of a series of
observations. A \emph{Type B evaluation} uses other information, such as a
calibration certificate, specification, previous data or a justified model of
a bounded influence.

These labels classify evaluation methods. They do not classify an effect as
random or systematic, and do not rank reliability. A calibration correction
can have a Type A evaluation; uncertainty in a fluctuating influence can have
a Type B evaluation \cite[2.3, 3.3.3--3.3.5]{JCGM1002008EMeasurement}.
\end{hbdefinition}

\begin{hbtheorem}{thm:p15-measurement-sample-covariance}{proved=annex:p15.measurement.sample-covariance}{Means and their covariance}
Let \(Q_1,\ldots,Q_n\), \(n\geq2\), be independent identically distributed
random vectors with mean \(\mu\) and finite covariance \(C\). Put
\[
 \bar Q=\frac1n\sum_{k=1}^n Q_k,\qquad
 S=\frac1{n-1}\sum_{k=1}^n(Q_k-\bar Q)(Q_k-\bar Q)^{\mathsf T}.
\]
Then \(E\bar Q=\mu\), \(\operatorname{Cov}(\bar Q)=C/n\), and \(ES=C\).
Thus \(S/n\) is an unbiased estimator of the covariance of the mean.
Its scalar diagonal gives
\[
 u^2(\bar q)=\frac1{n(n-1)}\sum_k(q_k-\bar q)^2.
\]
For arbitrary dependent observations, the covariance of their mean is
\(n^{-2}\sum_{k,l}\operatorname{Cov}(Q_k,Q_l)\), not automatically \(C/n\).
\end{hbtheorem}

The square root of an unbiased variance estimator need not be unbiased for a
standard deviation. The guide's \(s/\sqrt n\) is an estimated standard
uncertainty, not the realized error of the mean. Its \(n-1\) degrees of freedom
refer to the simple independent-sample variance calculation. A fitted curve,
pooled variance or time series requires the appropriate estimation model
\cite[4.2]{JCGM1002008EMeasurement}. If a stable pooled standard deviation
\(s_{\rm p}\) describes the same repeated-reading process, the guide uses
\(s_{\rm p}/\sqrt n\); its variance-estimation degrees of freedom come from the
pool, not the number of new readings.

\begin{hbexample}{def:p15-measurement-temperature-data}{The guide's twenty-temperature illustration}
The observations in degrees Celsius are
\[
\begin{gathered}
96.90,\ 98.18,\ 98.25,\ 98.61,\ 99.03,\\
99.49,\ 99.56,\ 99.74,\ 99.89,\ 100.07,\\
100.33,\ 100.42,\ 100.68,\ 100.95,\ 101.11,\\
101.20,\ 101.57,\ 101.84,\ 102.36,\ 102.72 .
\end{gathered}
\]
Their mean is \(100.145\), their sample standard deviation is about \(1.489\),
and the estimated standard uncertainty of the mean is about \(0.333\).
The source's illustrative generating model has mean \(100\) and standard
deviation \(1.5\). The histogram in Figure 1 groups the observations into
unit-width bins and distinguishes spread of readings from uncertainty of the
mean. These are illustrative data, not a reported new experiment
\cite[4.4.1--4.4.3, Table~1 and Fig.~1]{JCGM1002008EMeasurement}.
\end{hbexample}

\begin{hbtheorem}{thm:p15-measurement-bounded-distributions}{proved=annex:p15.measurement.bounded-distributions}{Variance under declared bounded-distribution models}
Let \(a>0\) and \(c\in\mathbb R\).
A uniform distribution on \([c-a,c+a]\) has mean \(c\) and variance \(a^2/3\).
The symmetric triangular density
\[
 p(t)=\frac{a-|t-c|}{a^2}\,\mathbf1_{\{|t-c|\leq a\}}
\]
has mean \(c\) and variance \(a^2/6\).
The symmetric trapezoid with base width \(2a\) and top width \(2a\beta\),
\(0\leq\beta\leq1\), has variance \(a^2(1+\beta^2)/6\).
Finally, if \(\Phi\) is uniform on \([0,2\pi]\), then
\(c+a\sin\Phi\) has mean \(c\) and variance \(a^2/2\).
\end{hbtheorem}

A support interval by itself does not entail a uniform distribution. Uniform,
triangular, trapezoidal and approximately normal assignments express different
information. The guide offers these as evaluation models
\cite[4.3.7--4.3.9]{JCGM1002008EMeasurement}. For the common interval
\([96,104]\) in degrees Celsius, its Figure 2 gives standard uncertainties
\(4/\sqrt3\approx2.31\) and \(4/\sqrt6\approx1.63\), respectively.
Neither is the \(0.333\) uncertainty of a mean obtained from twenty
observations.

For a density \(p\) relative to the chosen numerical coordinate, its
differential entropy is \(h(p)=-\int p(t)\log p(t)\,dt\), using natural
logarithms and the convention \(0\log0=0\). On a bounded interval the
entropy may equal \(-\infty\). The reference coordinate and constraints
are part of an entropy-maximization problem.

\begin{hbtheorem}{thm:p15-measurement-asymmetric-model}{proved=annex:p15.measurement.asymmetric-model}{Bounds with a separately specified mean}
Fix \(A<m<B\), \(b_-=m-A\) and \(b_+=B-m\).
Among densities on \([A,B]\) with mean \(m\), the maximum-differential-entropy
density is
\[
 q_\lambda(t)=\frac{e^{-\lambda(t-m)}}{Z(\lambda)},\qquad
 Z(\lambda)=\int_A^B e^{-\lambda(t-m)}\,dt,
\]
where the unique \(\lambda\in\mathbb R\) makes the mean \(m\).
Here \(T\) denotes the coordinate random variable with this density.
If \(m=(A+B)/2\), then \(\lambda=0\) and the density is uniform.
Otherwise
\[
 \operatorname{Var}_{q_\lambda}(T)
 =b_+b_- -\frac{b_+-b_-}{\lambda}.
\]
This is a mathematical consequence of the stated entropy criterion and
constraints, not a consequence of knowing the bounds alone.
\end{hbtheorem}

The guide's asymmetric example specifies \(16.52\), lower bound \(16.40\)
and upper bound \(16.92\), in units \(10^{-6}\,{}^\circ\mathrm C^{-1}\).
A uniform model on those bounds instead has mean \(16.66\).
Its standard deviation \(0.52/\sqrt{12}\approx0.15\) belongs to that
uniform model; retaining the different mean requires another distribution.
The preceding result supplies the guide's maximum-entropy alternative
\cite[4.3.8 and its notes]{JCGM1002008EMeasurement}.

Write \(\Phi\) for the standard normal cumulative distribution function
in the following conversions. Further source conversions illustrate what
their premises supply:
\begin{itemize}
\item A mass-certificate uncertainty \(240\,\mu\mathrm g\) stated at three
standard deviations gives standard uncertainty \(80\,\mu\mathrm g\).
\item An approximately normal 99-percent half-width \(129\,\mu\Omega\)
gives about \(50\,\mu\Omega\) after division by \(2.58\).
A coverage probability without a distribution or multiplier does not
determine a standard uncertainty.
\item A normal central 50-percent interval with half-width \(a\) gives
\(u=a/\Phi^{-1}(0.75)\approx1.48a\). Thus the guide's
\(10.11\pm0.04\,\mathrm{mm}\) illustration gives about \(0.06\,\mathrm{mm}\).
A roughly two-thirds normal interval gives \(u\) approximately equal to its
half-width; extra digits would overstate the supplied information.
\item For the symmetric copper-coefficient bound
\(16.52\pm0.40\) in units \(10^{-6}\,{}^\circ\mathrm C^{-1}\),
a uniform assignment gives \(u\approx0.23\) in the same units.
\item An independent voltmeter correction with uniform half-width
\(15\,\mu\mathrm V\) has variance \(75\,\mu\mathrm V^2\).
Combined with a \(12\,\mu\mathrm V\) standard uncertainty of the observed
mean, it gives \(\sqrt{144+75}\approx15\,\mu\mathrm V\) at
\(0.928571\,\mathrm V\).
\end{itemize}
These are conditional reconstructions of the guide's examples
\cite[4.3.3--4.3.9, 5.1.5]{JCGM1002008EMeasurement}.
A component already represented in observed variability must not also be
added as an independent uncertainty for the same effect.

\section{Linear propagation and shared influences}
\label{sec:p15-measurement-linear}

\begin{hbtheorem}{thm:p15-measurement-affine-covariance}{proved=annex:p15.measurement.affine-covariance}{Exact affine propagation}
For \(X\) as in \cref{def:p15-measurement-model}, a constant matrix \(J\)
and vector \(a\), let \(Y=a+JX\). Then
\[
 EY=a+Jm,\qquad \operatorname{Cov}(Y)=J\Sigma J^{\mathsf T}.
\]
In the scalar case with coefficient row \(c^{\mathsf T}\),
\[
 \operatorname{Var}(Y)=\sum_i c_i^2\Sigma_{ii}
       +2\sum_{i<j}c_i c_j\Sigma_{ij}.
\]
Uncorrelated inputs suffice to remove the cross terms in this identity;
independence is stronger. For an assessed covariance matrix \(U\), the
corresponding assessed affine variance is \(c^{\mathsf T}Uc\).
\end{hbtheorem}

For nonlinear \(f\), evaluate the sensitivity row at the chosen center:
\(c_i=\partial_i f(x)\). The guide's first-order combined variance is
\[
 u_{\rm lin}^2(y)=c^{\mathsf T}Uc
   =\sum_i c_i^2u_i^2+2\sum_{i<j}c_i c_j u_i u_j r_{ij}.
\]
The last representation uses only pairs with positive standard deviations.
For a vector of outputs, use \(JUJ^{\mathsf T}\), including off-diagonal
entries \cite[5.1--5.2, H.2]{JCGM1002008EMeasurement}.
A positive covariance can increase or decrease the output variance depending
on the signs of the sensitivities.

If all nonzero-variance inputs have correlation \(+1\), the scalar standard
deviation is \(\left|\sum_i c_i u_i\right|\). If they are uncorrelated it
is \((\sum_i c_i^2u_i^2)^{1/2}\). Neither prescription is the general rule
for an arbitrary dependence structure.

\begin{hbexample}{def:p15-measurement-shared-reference}{A shared calibration influence}
Let \(A=a+C+E_1\), \(B=b+C+E_2\), with independent centered \(C,E_1,E_2\)
and variances \(\tau^2,s_1^2,s_2^2\). Then
\[
 \operatorname{Var}(A-B)=s_1^2+s_2^2,\qquad
 \operatorname{Var}(A+B)=4\tau^2+s_1^2+s_2^2.
\]
Dropping the shared covariance gives a wrong answer in both cases.
The guide's ten-resistor example is a sum with one common calibration
uncertainty: ten \(0.1\,\Omega\) contributions with correlation \(+1\)
give \(1\,\Omega\), not \(\sqrt{10}\,0.1\,\Omega\)
\cite[5.2.2]{JCGM1002008EMeasurement}.
\end{hbexample}

\begin{hbtheorem}{thm:p15-measurement-common-correction}{proved=annex:p15.measurement.common-correction}{Averaging does not remove a common uncertain correction}
Let \(A,B,Q\) be independent with finite second moments, means \(a,b,q\)
and variances \(v_A,v_B,v_Q\). For \(W=A+BQ\),
\[
 EW=a+bq,\qquad
 \operatorname{Var}(W)=v_A+q^2v_B+b^2v_Q+v_Bv_Q.
\]
The first-order expression omits \(v_Bv_Q\). If \(Q\) is the mean of
\(n\) independent readings with variance \(\sigma^2\), then
\(v_Q=\sigma^2/n\), but the first two variance terms need not shrink.
\end{hbtheorem}

This makes the distinction in Annex E.3.4--E.3.7 explicit: a common offset and
scale can be fixed during data collection yet uncertain in the measurement
model. Treating the offset as an independent new draw for each reading changes
the model. The identical variance algebra for Type A and Type B inputs is
conditional on their joint probabilistic interpretation.

Annex E.3.5 illustrates the degree-of-belief interpretation by comparing a
payment \(D\) on an event with the same payment on its complement. Under
expected utility, the same two consequences, and a strict preference for
receiving \(D\), indifference between these bets gives \(P(A)=1/2\).
Those decision assumptions explain the inference; indifference by itself
does not prove a probability representation. The relevant representation
questions belong to the separate subjective-probability treatments.

For a product \(f(x)=K\prod_i x_i^{p_i}\) on a differentiable nonzero domain,
\(c_i=f(x)p_i/x_i\). Thus the relative first-order variance has diagonal
terms \(p_i^2u_i^2/x_i^2\) and covariance terms
\(2p_i p_j u_{ij}/(x_i x_j)\). Zero inputs or outputs cannot be handled
by dividing by them. Logarithmic linearization additionally requires positive
arguments \cite[5.1.6, 5.2.2]{JCGM1002008EMeasurement}.
Sensitivity coefficients may instead be measured or approximated numerically;
finite perturbations require their own approximation check.

\section{What nonlinear propagation requires}
\label{sec:p15-measurement-nonlinear}

\begin{hbtheorem}{thm:p15-measurement-taylor-bound}{proved=annex:p15.measurement.taylor-bound}{A controlled first-order approximation}
Use fixed numerical coordinates. Suppose \(f\) is twice continuously
differentiable on an open convex set containing \(m\) and \(X\) almost surely,
and its Hessian has operator norm at most \(M\) there. Assume
\(E\|X-m\|^4<\infty\). Use the Euclidean vector norm and its induced
matrix norm, and \(\|R\|_2=(E|R|^2)^{1/2}\) for a scalar random remainder. Put
\[
 L=f(m)+Df(m)(X-m),\quad R=f(X)-L,\quad
 r=\tfrac M2\{E\|X-m\|^4\}^{1/2}.
\]
Then
\[
 |Ef(X)-f(m)|\leq \tfrac M2 E\|X-m\|^2,\qquad
 \|R\|_2\leq r,
\]
and
\[
 \left|\operatorname{Var}(f(X))-\operatorname{Var}(L)\right|
 \leq2\sqrt{\operatorname{Var}(L)}\,r+r^2.
\]
These are bounds under the displayed assumptions, not claims about arbitrary
measurement functions.
\end{hbtheorem}

The output mean and the output variance require separate checks.
Even when \(x=m\), \(f(x)\) is not generally \(Ef(X)\). For a quadratic
function with symmetric Hessian \(H\), the exact mean is
\[
 Ef(X)=f(m)+\tfrac12\sum_{i,j}H_{ij}\Sigma_{ij}.
\]
This identity needs second moments, not normality.
The amendment adds a corresponding nonlinear mean qualification to 4.1.4,
with the independent-normal diagonal expression and a reference to H.10
\cite{JCGM100Amendment2026Nonlinearity}.

\begin{hbtheorem}{thm:p15-measurement-gaussian-expansion}{proved=annex:p15.measurement.gaussian-expansion}{The Gaussian next-order terms}
Let \(Z_i\) be independent centered normal variables with variances
\(\sigma_i^2\geq0\), for \(1\leq i\leq d<\infty\).
All indices in this result range from \(1\) to \(d\).
Fix coefficients \(g_i\), a symmetric matrix \(H\)
and a fully symmetric array \(T_{ijk}\). Define
\[
 A=\sum_i g_iZ_i,\quad B=\tfrac12\sum_{i,j}H_{ij}Z_iZ_j,\quad
 C=\tfrac16\sum_{i,j,k}T_{ijk}Z_iZ_jZ_k,\quad
 F_\varepsilon=f_0+\varepsilon A+\varepsilon^2B+\varepsilon^3C.
\]
Then
\[
 EF_\varepsilon=f_0+\tfrac{\varepsilon^2}2\sum_iH_{ii}\sigma_i^2
\]
and
\[
\begin{split}
 \operatorname{Var}(F_\varepsilon)
  ={}&\varepsilon^2\sum_i g_i^2\sigma_i^2\\
    &+\varepsilon^4\sum_{i,j}
       \left(\tfrac12H_{ij}^2+g_iT_{ijj}\right)\sigma_i^2\sigma_j^2
      +\varepsilon^6\operatorname{Var}(C).
\end{split}
\]
For a cubic polynomial this is an exact calculation.
For a smooth nonpolynomial function these are coefficients of a Taylor
surrogate; applying them to the function requires control of the omitted
remainder.
\end{hbtheorem}

The gradient--third-derivative term must accompany the squared-Hessian
term at this variance order. Retaining a quadratic surrogate alone does not
in general supply the full next-order variance of a nonlinear function.
The source's note to 5.1.2 is read under its independent-input, normal-model
conditions \cite[5.1.2]{JCGM1002008EMeasurement}.

\begin{hbexample}{def:p15-measurement-zero-gradient}{A zero gradient with a nonzero uncertainty}
If \(X\) is uniform on \([-a,a]\), \(a>0\), then
\[
 EX=0,\quad E(X^2)=a^2/3,\quad
 \operatorname{Var}(X^2)=a^4/5-a^4/9=4a^4/45.
\]
Linearizing \(x^2\) at zero gives both output estimate and propagated
variance zero. Both conclusions miss the stipulated distribution.
The moments follow by integrating \(x^2\) and \(x^4\) against \(1/(2a)\).
Nondegeneracy of \(X\) alone would not imply nondegeneracy of \(X^2\):
a symmetric two-point law on \(\{-a,a\}\) has constant square.
\end{hbexample}

\section{Coverage factors and what an interval means}
\label{sec:p15-measurement-coverage}

\begin{hbdefinition}{def:p15-measurement-expanded}{Expanded uncertainty}
Given an output estimate \(y\), a combined standard uncertainty \(u_c\)
and a stated \(k>0\), the expanded uncertainty is \(U=ku_c\).
The interval is \([y-U,y+U]\). Assigning it a probability requires a
distributional interpretation. A fixed-parameter confidence claim instead
requires the repeated-sampling coverage of an interval-producing procedure.
\end{hbdefinition}

\begin{hbtheorem}{thm:p15-measurement-coverage-models}{proved=annex:p15.measurement.coverage-models}{A multiplier is not a universal coverage probability}
Let \(k>0\), \(a>0\) and \(c\in\mathbb R\).
For a centered normal output with known standard deviation \(\sigma>0\),
the interval of radius \(k\sigma\) around its mean has probability
\(2\Phi(k)-1\). For a uniform output on \([c-a,c+a]\), using
\(\sigma=a/\sqrt3\), that probability is
\(\min\{k/\sqrt3,1\}\).
For the symmetric triangular model, using \(\sigma=a/\sqrt6\), it is
\(2z-z^2\) for \(z=k/\sqrt6\leq1\), and one otherwise.
For any output of finite variance \(\sigma^2>0\), Chebyshev gives only
\[
 P(|Y-EY|<k\sigma)\geq1-k^{-2}\quad(k>1).
\]
\end{hbtheorem}

The usual normal \(k=2\) probability is about \(0.9545\); \(k=3\) gives
about \(0.9973\). A numerical multiplier does not create additional
information. Estimated scale, non-normality and tail uncertainty matter
\cite[6.2--6.3]{JCGM1002008EMeasurement}.
The exact Student confidence procedure for independent normal readings
is already given by \cref{thm:p15-student-law,thm:p15-student-coverage}.
Its sampling guarantee is not a general property of a mixed Type A/Type B
measurement budget.

\begin{hbtheorem}{thm:p15-measurement-effective-df}{proved=annex:p15.measurement.effective-df}{What the effective-degrees-of-freedom formula matches}
Let independent nonnegative estimators \(V_i\) have
\(EV_i=v_i\geq0\), \(\operatorname{Var}(V_i)=2v_i^2/\nu_i\),
where \(\nu_i>0\), and let \(V=\sum_i V_i\), \(v=\sum_i v_i>0\).
A scaled chi-square variable \(v\,\chi_\nu^2/\nu\) matches the mean and
variance of \(V\) when
\[
 \nu=\frac{v^2}{\sum_i v_i^2/\nu_i}.
\]
This matching does not assert equality of distributions.
\end{hbtheorem}

The guide's Welch--Satterthwaite prescription substitutes assessed
contributions \(v_i=c_i^2u_i^2\), obtains
\(\nu_{\rm eff}=u_c^4/\sum_i(c_i u_i)^4/\nu_i\), and uses the corresponding
Student quantile as an approximation \cite[G.4]{JCGM1002008EMeasurement}.
It does not automatically apply to correlated components or an incomplete
nonlinear budget. The guide's Type B reliability approximation
\(\nu_i\approx[2(\Delta u_i/u_i)^2]^{-1}\) assigns \(8,50,2\) to relative
reliabilities \(25,10,50\) percent. This is a stated assessment convention,
not observed degrees of freedom created by those percentages.

\section{The end-gauge calculation with the nonlinear terms included}
\label{sec:p15-measurement-gauge}

The target is length at \(20\,{}^\circ\mathrm C\), compared with a standard
of the same nominal \(50\,\mathrm{mm}\) length. The source's physical model is
\[
 d=l(1+\alpha\theta)-l_{\rm S}(1+\alpha_{\rm S}\theta_{\rm S}),
 \qquad
 l=\frac{l_{\rm S}(1+\alpha_{\rm S}\theta_{\rm S})+d}{1+\alpha\theta}.
\]
Write \(A=\alpha-\alpha_{\rm S}\), \(D=\theta-\theta_{\rm S}\).
The source's approximation uses
\[
 l\approx l_{\rm S}+d-l_{\rm S}(A\theta+\alpha_{\rm S}D).
\]
The exact omitted term relative to this approximation is
\[
 \frac{\alpha\theta\{l_{\rm S}(A\theta+\alpha_{\rm S}D)-d\}}
      {1+\alpha\theta}.
\]
Thus small thermal products and an admissible denominator are assumptions
of the simplified model, rather than exact cancellations
\cite[H.1.1--H.1.2]{JCGM1002008EMeasurement}.

For the uncertainty budget below take the nominal \(L=0.05\,\mathrm m\)
as the sensitivity length. The guide's input means are
\(EA=ED=0\), \(E\theta=-0.1\,{}^\circ\mathrm C\) and
\(E\alpha_{\rm S}=11.5\,10^{-6}\,{}^\circ\mathrm C^{-1}\).
The assigned standard uncertainties are:
\begin{center}
\begin{tabularx}{\linewidth}{@{}Xll@{}}
\toprule
Input or component & Evaluation & Standard uncertainty\\
\midrule
Standard length & \(0.075\,\mu\mathrm m/3\) & \(25\,\mathrm{nm}\)\\
Repeated comparison mean & \(13\,\mathrm{nm}/\sqrt5\) & \(5.8\,\mathrm{nm}\)\\
Comparator random-effect correction & \(10\,\mathrm{nm}/2.57\) & \(3.9\,\mathrm{nm}\)\\
Comparator systematic-effect correction & \(20\,\mathrm{nm}/3\) & \(6.7\,\mathrm{nm}\)\\
Standard expansion coefficient & \(2\,10^{-6}/\sqrt3\) & \(1.2\,10^{-6}\,{}^\circ\mathrm C^{-1}\)\\
Mean bed temperature & supplied & \(0.2\,{}^\circ\mathrm C\)\\
Cyclic temperature component & \(0.5/\sqrt2\) & \(0.35\,{}^\circ\mathrm C\)\\
Expansion-coefficient difference \(A\) & \(10^{-6}/\sqrt3\) & \(0.58\,10^{-6}\,{}^\circ\mathrm C^{-1}\)\\
Temperature difference \(D\) & \(0.05/\sqrt3\) & \(0.029\,{}^\circ\mathrm C\)\\
\bottomrule
\end{tabularx}
\end{center}
The comparison contributions give \(u^2(d)\approx93.4\,\mathrm{nm}^2\);
the temperature contributions give \(u^2(\theta)=0.165\,{}^\circ\mathrm C^2\).
The \(0.5\)-degree cyclic amplitude is not an uncertainty of the mean
temperature: uniform phase produces the sine-model variance already derived
\cite[H.1.3]{JCGM1002008EMeasurement}.

\begin{hbtheorem}{thm:p15-measurement-gauge-products}{proved=annex:p15.measurement.gauge-products}{The bilinear thermal contribution}
Let \(A,\Theta,B,D\) be mutually independent with finite second moments,
\(EA=ED=0\), \(E\Theta=\theta\), \(EB=b\), and let \(E_0\) be an
independent centered error of variance \(v_0\). For
\[
 G=g_0+E_0-L(A\Theta+BD),
\]
the mean is \(g_0\) and the variance is exactly
\[
 v_0+L^2\{\theta^2v_A+b^2v_D+v_Av_\Theta+v_Bv_D\}.
\]
Here \(v_A,v_\Theta,v_B,v_D\) denote the four variances.
\end{hbtheorem}

With the source inputs, the linearized part is approximately
\(1002\,\mathrm{nm}^2\). The two product contributions are approximately
\(137.5\,\mathrm{nm}^2\) and \(2.78\,\mathrm{nm}^2\).
The resulting standard uncertainty is approximately \(33.8\,\mathrm{nm}\),
reported as \(34\,\mathrm{nm}\). This moment calculation makes the
independence assumptions explicit; mere pairwise uncorrelatedness would not
determine product variances. The guide's H.1.7 uses the same two
next-order contributions, and the amendment emphasizes checking the mean
as well as the variance. The mean correction vanishes under the centered,
uncorrelated thermal-difference assumptions
\cite[H.1.7]{JCGM1002008EMeasurement}
\cite{JCGM100Amendment2026Nonlinearity}.

Using \(l_{\rm S}=50.000623\,\mathrm{mm}\) and mean difference
\(215\,\mathrm{nm}\), the result within the specified approximation is
\[
 l=50.000838\,\mathrm{mm},\qquad u_c(l)\approx34\,\mathrm{nm}.
\]
Its relative standard uncertainty is about \(6.8\,10^{-7}\).
Keeping uncertainty in the length multiplying the thermal term adds an
additional higher-order product contribution; its neglect, like the
rational-model remainder above, belongs to the approximation budget.

The source's linear budget illustrates effective degrees of freedom:
the comparison pool supplies \(24\), the comparator random component \(5\),
and the assessed systematic component \(8\), leading to about \(25.6\)
for their combined contribution. The other listed degrees are \(18\) for
the standard and assessed values \(50,2\) for \(A,D\).
Applying the prescription to that \emph{linear} budget gives the source's
approximately \(16.7\), truncated to \(16\), and \(t_{16,0.995}\approx2.92\).
These numbers describe that particular approximate budget.
They do not establish a 99-percent interval for the completed nonlinear
calculation: reliability of the added product terms and the adequacy of the
output distribution must also be assessed. The standard-uncertainty result
above therefore does not silently inherit a coverage probability
\cite[H.1.6--H.1.7]{JCGM1002008EMeasurement}.

\section{Correlated electrical inputs and three related outputs}
\label{sec:p15-measurement-impedance}

For voltage amplitude \(V\geq0\), current amplitude \(I>0\) and phase
\(\phi\) in radians, the source defines resistance \(R\), reactance
\(X_{\rm r}\) (called \(X\) in H.2), and impedance magnitude \(Z\):
\[
 R=(V/I)\cos\phi,\qquad X_{\rm r}=(V/I)\sin\phi,\qquad Z=V/I.
\]
The five simultaneous observation sets are
\begin{center}
\begin{tabular}{rrrr}
\toprule
Set & \(V\;(\mathrm V)\) & \(I\;(\mathrm{mA})\) & \(\phi\;(\mathrm{rad})\)\\
\midrule
1&5.007&19.663&1.0456\\
2&4.994&19.639&1.0438\\
3&5.005&19.640&1.0468\\
4&4.990&19.685&1.0428\\
5&4.999&19.678&1.0433\\
\bottomrule
\end{tabular}
\end{center}
They give means \(4.9990\,\mathrm V\), \(19.6610\,\mathrm{mA}\),
\(1.04446\,\mathrm{rad}\), with standard uncertainties of the means
approximately \(0.003209\,\mathrm V\), \(0.009471\,\mathrm{mA}\) and
\(0.0007521\,\mathrm{rad}\). The three estimated input correlations are
approximately \(-0.3553,0.8576,-0.6451\), in the order
\((V,I),(V,\phi),(I,\phi)\)
\cite[H.2.1--H.2.2, Table~H.2]{JCGM1002008EMeasurement}.

\begin{hbtheorem}{thm:p15-measurement-impedance-jacobian}{proved=annex:p15.measurement.impedance-jacobian}{Propagation for the impedance coordinates}
The Jacobian with input coordinates \((V,I,\phi)\) is
\[
 J=\begin{pmatrix}
 \cos\phi/I&-V\cos\phi/I^2&-X_{\rm r}\\
 \sin\phi/I&-V\sin\phi/I^2&R\\
 1/I&-V/I^2&0
 \end{pmatrix}.
\]
The first-order output covariance is \(JUJ^{\mathsf T}\).
At \(V>0\), its magnitude coordinate satisfies
\[
 \frac{u_{\rm lin}^2(Z)}{Z^2}
 =\frac{u_V^2}{V^2}+\frac{u_I^2}{I^2}
       -\frac{2u_{VI}}{VI}.
\]
The identity \(Z^2=R^2+X_{\rm r}^2\) makes the three outputs dependent;
it does not make any two statistical errors independent.
\end{hbtheorem}

Recomputing from the five unrounded rows, with current converted to amperes,
gives:
\begin{center}
\begin{tabularx}{\linewidth}{@{}Xrrr@{}}
\toprule
 & \(R\;(\Omega)\)&\(X_{\rm r}\;(\Omega)\)&\(Z\;(\Omega)\)\\
\midrule
Transform of input means &127.732170&219.846512&254.259702\\
First-order standard uncertainty &0.071071&0.295582&0.236336\\
Mean of transformed rows &127.731630&219.846895&254.260050\\
Standard uncertainty of that mean &0.071274&0.295489&0.236248\\
First-order result after dropping input covariances &0.194544&0.200909&0.204076\\
\bottomrule
\end{tabularx}
\end{center}
Extra digits expose the comparison; they are not recommended reporting
precision. The first-order output correlations are approximately
\(-0.5884,-0.4853,0.9925\).
The source's rounded Tables H.3--H.5 give the corresponding small differences
and the substantial effects of dropping covariance
\cite[H.2.3--H.2.4]{JCGM1002008EMeasurement}.
The dataset accounts only for observed random variation; an actual circuit
measurement also needs uncertainty in significant systematic corrections.

\begin{hbtheorem}{thm:p15-measurement-transform-average}{proved=annex:p15.measurement.transform-average}{Transforming an average and averaging transforms}
For fixed vectors \(q_1,\ldots,q_n\), \(n\geq2\), let
\[
 \bar q=n^{-1}\sum_kq_k,\qquad
 U_{\rm mean}=\frac1{n(n-1)}
       \sum_k(q_k-\bar q)(q_k-\bar q)^{\mathsf T}.
\]
For affine \(f\), \(n^{-1}\sum_k f(q_k)=f(\bar q)\).
For quadratic \(f\) with Hessian \(H\), the exact difference is
\[
 \frac1n\sum_kf(q_k)-f(\bar q)
       =\frac{n-1}{2}\sum_{i,j}H_{ij}(U_{\rm mean})_{ij}.
\]
For \(f\) twice continuously differentiable on an open convex set containing
these points, the same identity uses the
appropriate integral Taylor remainders instead of one fixed Hessian.
\end{hbtheorem}

This finite-sample identity differs from
\(Ef(X)-f(EX)=\tfrac12\sum H_{ij}\Sigma_{ij}\) for a quadratic uncertainty
model: the latter covariance belongs to the random input being transformed.
The former covariance is estimated for a sample mean. The distinction matters
when using nonlinear-mean corrections.

Pairing must reflect data collection. The guide prefers transforming the
simultaneous triples in this example because that follows the observation
procedure. Five voltages observed in one period, five currents later and five
phases later do not supply the same pairing or empirical cross-covariance.
Arbitrary row matching does not reconstruct missing joint observations
\cite[4.1.4, H.2.4]{JCGM1002008EMeasurement}.

\section{A reference procedure and its stopping conditions}
\label{sec:p15-measurement-procedure}

\begin{hbmethod}{def:p15-measurement-procedure}{From a specified quantity to a defensible uncertainty report}
\begin{enumerate}
\item Define the measurand, units, reference conditions, realization and
intended comparison. Resolve a material ambiguity before calculating.
\item State the measurement equation or algorithm, including significant
corrections and influence quantities. Identify its validity domain.
\item Record input estimates, standard uncertainties, information sources,
evaluation methods and the meaning of any reliability or degrees of freedom.
\item Record covariance or a justified joint model. Identify shared
references and avoid counting the same effect twice.
\item Calculate the output estimate and the linearized covariance.
Check nonlinear effects on both its center and spread; include justified
higher-order terms or use a separately validated propagation procedure.
\item If an expanded interval is needed, state its multiplier, probability
interpretation and distributional or sampling justification. A tolerance,
safety allowance or known uncorrected bias is not an uncertainty component
merely because it is written after a plus-or-minus sign.
\item Report \(y,u_c\), their units, and relative uncertainty only if
\(y\ne0\). For several outputs report their covariance or correlation.
For expanded uncertainty also report \(k\), the intended probability and how
it was assessed. Preserve sufficient intermediate precision for reproduction.
\item State remaining limitations that could affect the inference, including
model incompleteness, missing covariance or an unsupported coverage claim.
If the proposed conclusion changes within the justified uncertainty range,
do not assert that conclusion from the point estimate alone.
\end{enumerate}
This reconstruction implements the source's reporting and calculation
recommendations with the explicit nonlinear distinctions above
\cite[7--8]{JCGM1002008EMeasurement}.
\end{hbmethod}

The guide's reporting examples distinguish a standard uncertainty
\(0.35\,\mathrm{mg}\) attached to \(100.02147\,\mathrm g\) from an expanded
uncertainty \(0.79\,\mathrm{mg}\) obtained with \(k=2.26\) and a stated
Student approximation. A parenthetic last-digits notation, a separately
named uncertainty or a plus-or-minus notation must make clear which quantity
is reported. Typically one or two uncertainty digits suffice for a final
report, while calculations can retain more \cite[7.2]{JCGM1002008EMeasurement}.

Changing the reference can also change the measurand: a ratio to an adopted
reference value and a result expressed in physical units have different
uncertainty budgets. The guide's historical voltage example compares
relative uncertainties \(2\,10^{-8}\) and \(4\,10^{-7}\) for those two
reference choices; the numerical example belongs to its stated historical
reference conventions \cite[3.4.6]{JCGM1002008EMeasurement}.

A complete budget can be mathematically consistent yet poorly supported
empirically. Critical examination of the measurement model, corrections and
information sources remains part of the inference. The measurement result
is usable when its conditions match the next calculation, not merely because
an uncertainty value has been printed.
